\chapter{Cài đặt và tích hợp hệ thống}
\markboth{Chương 4. Cài đặt hệ thống}{}
\section{Công nghệ và tổ chức mã nguồn}
Frontend sử dụng React 19, Vite 8, Tailwind CSS 4, MediaPipe Tasks Vision, Framer Motion và Lucide. Backend dùng FastAPI, Uvicorn, SQLAlchemy, APScheduler, python-dotenv, openpyxl, python-docx và python-multipart. Flutter nằm trong thư mục mobile, là phần mở rộng độc lập với frontend web.

Các phụ thuộc JavaScript có package-lock.json; requirements.txt của Python chủ yếu khai báo phiên bản tối thiểu. Vì vậy, môi trường kiểm thử cần lưu phiên bản thực tế đã cài để tái lập. Mô hình được tải từ đường dẫn phiên bản cụ thể, nhưng WASM hiện sử dụng đường dẫn latest; đây là giới hạn đối với tính ổn định khi nguồn ngoài thay đổi.

\section{Tài khoản và phiên truy cập}
\subsection{Giảng viên}
Tài khoản giảng viên được tạo bằng create\_teacher.py. Mã theo định dạng GV-xxxx; mật khẩu được nhập tương tác. auth.py dùng scrypt cùng salt ngẫu nhiên; đường kiểm tra còn hỗ trợ hash PBKDF2 cũ và nâng hash sau đăng nhập thành công. CSDL không lưu mật khẩu dạng rõ.

Đăng nhập thành công tạo token ngẫu nhiên với thời hạn 12 giờ. Máy chủ lưu SHA-256 của token trong teacher\_sessions; frontend giữ token trong sessionStorage và gửi qua Authorization: Bearer. Đăng xuất vô hiệu hóa phiên phía server. Có giới hạn số lần đăng nhập sai theo IP. Cơ chế token legacy chỉ được bật khi cấu hình ALLOW\_LEGACY\_TOKEN=1 và có bí mật tương ứng; đây không phải luồng đăng nhập chính.

\subsection{Sinh viên và phê duyệt}
Mã sinh viên mặc định có dạng SV- theo sau bởi 8 đến 12 chữ số, có thể thay đổi bằng STUDENT\_ID\_REGEX. Sau khi gửi yêu cầu vào thi, sinh viên nhận request token và dùng header X-Request-Token để theo dõi trạng thái. Phê duyệt chưa trực tiếp trả token thi cho giảng viên; lần kiểm tra trạng thái hợp lệ sau đó cấp token thi một lần cho sinh viên. Token được truyền bằng X-Exam-Session ở các yêu cầu lấy đề, bắt đầu, nộp và gửi vi phạm.

Dependency get\_exam\_session kiểm tra bản băm token, trạng thái APPROVED và cờ đã cấp phiên. Backend lấy student\_id, student\_name và exam\_id từ lượt ghi danh. Các schema bài nộp và vi phạm không chấp nhận trường ngoài quy định, giúp tránh việc client tự thay danh tính trong những payload này.

\section{Nhập đề và quản lý lịch}
API upload nhận tên đề, môn học, mã tùy chọn và số phút cho phép xin vào sớm. Tệp giới hạn 5 MiB; bộ đề tối đa 200 câu hợp lệ. Backend đọc worksheet đầu của Excel theo bảy cột, hoặc các đoạn Word theo mẫu câu hỏi--A/B/C/D--đáp án. Dòng thiếu dữ liệu, đáp án sai định dạng và câu trùng được ghi nhận khi phân tích.

Câu trùng được xác định bằng nội dung đã strip và chuyển chữ thường, không phải so khớp ngữ nghĩa. Sau xử lý, thứ tự câu được đánh lại và trọng số mặc định chia đều theo thang 10. Mã đề được chuẩn hóa về ký tự cho phép và kiểm tra duy nhất. Đề mới gắn với tài khoản giảng viên đang đăng nhập.

Công bố yêu cầu thời lượng 1--600 phút, ngưỡng đạt 0--10 và hai mốc mở/đóng hợp lệ. Giao diện cho phép nhập lịch và chuyển thời gian sang dạng gửi API. Sinh viên chỉ truy cập theo điều kiện lịch và phiên của đề. Xáo trộn áp dụng cho thứ tự câu hỏi, không phải hoán đổi các phương án A--D. Thứ tự đã hiển thị có thể được giữ ở tiến độ cục bộ trong phiên trình duyệt.

\section{Mốc bắt đầu, khôi phục và hạn nộp}
Lần bắt đầu đầu tiên tạo exam\_starts với ràng buộc duy nhất theo đề và sinh viên. Gọi lại sử dụng mốc đã có để tính thời gian còn lại; không tạo thời lượng mới. Frontend xác định thời điểm kết thúc cục bộ bằng thời gian còn lại từ server và dùng Date.now() khi cập nhật đồng hồ, thay vì chỉ trừ dần từng giây.

sessionStorage lưu đáp án theo question.id, câu hiện tại, thứ tự câu và thông tin phiên. Khi tải lại, frontend lọc các đáp án còn hợp lệ, phục hồi thứ tự nếu khớp bộ câu và gọi lại API để đồng bộ thời gian. Cơ chế này hỗ trợ tải lại trong phiên trình duyệt, chưa phải lưu nháp trên server hoặc khôi phục từ thiết bị khác.

Đặt $t_s$ là giờ bắt đầu tại server, $D$ là thời lượng tính bằng giây và $t_e$ là giờ đóng đề. Hạn tiếp nhận bài nộp là:
\begin{equation}
 t_{\mathrm{deadline}}=\min(t_s+D+15\text{ giây},\ t_e).
\end{equation}
Khoảng 15 giây dành cho độ trễ gửi bài, không kéo dài giờ đóng đề. API từ chối khi thiếu mốc bắt đầu hoặc giờ server vượt hạn. Cờ auto\_submitted được dùng ghi trạng thái nộp, không thay thế kiểm tra thời gian tại backend.

\section{Chấm bài và xuất kết quả}
Backend truy vấn đáp án chuẩn, tạo ánh xạ lựa chọn theo ID câu và cộng trọng số câu đúng. Điểm được giới hạn tối đa 10 và làm tròn hai chữ số. Bài nộp cùng các câu trả lời được ghi trong giao dịch. Ràng buộc UNIQUE theo đề và sinh viên là lớp cuối chống bản nộp trùng; xung đột được trả về HTTP 409.

Bảng điểm gồm mã, tên sinh viên, điểm, số câu đúng, thời điểm nộp, trạng thái và số vi phạm. Phép đếm vi phạm dùng student\_id cùng exam\_id, không ghép đơn thuần theo khoảng thời gian. Giảng viên có thể xuất bảng điểm Excel. Số vi phạm không được đưa vào công thức trừ điểm.

\section{Khởi tạo camera và mô hình AI}
Camera được yêu cầu với kích thước mong muốn 640$\times$480 và hướng người dùng; độ phân giải thực tế phụ thuộc thiết bị. Face Landmarker chạy VIDEO, tối đa ba khuôn mặt, ngưỡng phát hiện và theo dõi 0,5; ưu tiên GPU rồi thử CPU khi không khởi tạo được. EfficientDet Lite0 chạy CPU, tối đa năm kết quả và ngưỡng 0,25.

Nếu mô hình vật thể tải thất bại, code tiếp tục với phần khuôn mặt và ghi cảnh báo; tính năng phát hiện điện thoại khi đó không hoạt động đầy đủ. Các mô hình được dùng để suy luận từ trọng số có sẵn, không có pipeline huấn luyện trong dự án.

\begin{table}[htbp]\centering
\caption{Thông số AI và ảnh sự kiện trên web}
\begin{tabularx}{\textwidth}{@{}L{6cm}X@{}}\toprule
Thông số & Giá trị \\\midrule
Chu kỳ phân tích & PC 300 ms; điện thoại 500 ms.\\
Khoảng nghỉ cùng nhãn & 5 giây.\\
Điểm điện thoại & Nhãn cell phone, score $>0{,}25$.\\
Landmarks dùng quay đầu & Mũi 1; mắt 133 và 362.\\
Ngưỡng hình học & Độ lệch $>0{,}25$ hoặc bất đối xứng $>1{,}25$.\\
Quay đầu liên tiếp & 12 lượt PC; 8 lượt mobile.\\
Độ sáng & Canvas 64$\times$48; trung bình $<15$.\\
Ảnh bằng chứng & JPEG, rộng 320 px, giữ tỷ lệ, chất lượng 0,6.\\
\bottomrule\end{tabularx}\end{table}

\section{Quy tắc chọn nhãn và gửi sự kiện}
\begin{algorithm}[H]
\caption{Chọn trạng thái giám sát từ một lượt phân tích}
\begin{minipage}{0.94\linewidth}
\small
\textbf{Đầu vào:} khung hình, kết quả mặt/vật thể, bộ đếm $q$, chu kỳ $\Delta t$.\\
\textbf{Đầu ra:} nhãn chính $s$, bộ đếm mới.
\begin{enumerate}[leftmargin=1.6em,itemsep=2pt]
\item Lưu $q_{old}\gets q$ và đặt $q\gets0$.
\item Nếu độ sáng $<15$: chọn CAMERA\_BLOCKED.
\item Ngược lại, nếu có cell phone với score $>0{,}25$: chọn PHONE\_DETECTED.
\item Ngược lại, nếu nhiều hơn một mặt: chọn MULTIPLE\_FACES; nếu không có mặt: chọn NO\_FACE.
\item Với đúng một mặt, tính độ lệch $r$ và bất đối xứng $a$. Nếu $r>0{,}25$ hoặc $a>1{,}25$, đặt $q\gets q_{old}+1$.
\item Riêng nhánh đúng một mặt: chọn TURNING\_HEAD khi $q\geq\lceil3600/\Delta t\rceil$; nếu chưa đủ, chọn NORMAL.
\item Nếu nhãn khác NORMAL và đủ cooldown, gửi sự kiện cùng ảnh.
\end{enumerate}
\end{minipage}
\end{algorithm}

Mỗi lượt chỉ chọn một nhãn theo ưu tiên. Khi điện thoại và nhiều khuôn mặt cùng xuất hiện, nhãn điện thoại được chọn trước. Số bản ghi còn phụ thuộc cooldown, vì vậy không đồng nhất số bản ghi với tổng số hành vi hoặc tổng thời gian vi phạm.
\begin{figure}[H]\centering
\begin{tikzpicture}[x=1cm,y=1cm, every node/.style={font=\fontsize{10}{12}\selectfont}]
\node[rb,font=\fontsize{10}{12}\selectfont] (cam) at (0,0) {Khung hình camera\\300 / 500 ms};
\node[rb,font=\fontsize{10}{12}\selectfont] (face) at (-2.8,-2.4) {Face Landmarker\\Số mặt / landmarks};
\node[rb,font=\fontsize{10}{12}\selectfont] (obj) at (2.8,-2.4) {EfficientDet + Canvas\\Điện thoại / độ sáng};
\node[rb,text width=9cm,font=\fontsize{10}{12}\selectfont] (rule) at (0,-4.8) {Ưu tiên: tối → điện thoại → nhiều mặt\\→ không mặt → quay đầu → bình thường};
\node[rb,font=\fontsize{10}{12}\selectfont] (normal) at (-2.8,-7.2) {NORMAL\\Tiếp tục quan sát};
\node[rb,font=\fontsize{10}{12}\selectfont] (bad) at (2.8,-7.2) {Bất thường\\Kiểm tra cooldown 5 s};
\node[rb,font=\fontsize{10}{12}\selectfont] (send) at (2.8,-9.6) {Ảnh + metadata\\POST với phiên thi};
\draw[ra] (cam)--(face);\draw[ra] (cam)--(obj);\draw[ra] (face)--(rule);\draw[ra] (obj)--(rule);
\draw[ra] (rule)--(normal);\draw[ra] (rule)--(bad);\draw[ra] (bad)--(send);
\end{tikzpicture}

\caption{Luồng xử lý hình ảnh và gửi cảnh báo phía sinh viên.}\end{figure}

Sự kiện trình duyệt đi qua callback riêng. TAB\_SWITCH, WINDOW\_BLUR và FULLSCREEN\_EXIT được gửi khi bài thi hoạt động. Hook gộp các lần rời trang gần nhau trong 1,5 giây; bước gửi áp dụng cooldown theo nhãn. Các thao tác như phím bị chặn hoặc chuột phải có thể chỉ tạo thông báo cục bộ. JavaScript không cung cấp khóa hệ điều hành.

\section{Tiếp nhận sự kiện và bảo vệ lưu trữ}
Backend kiểm tra phiên, loại vi phạm và thời điểm chụp lệch không quá 300 giây so với server. Tần suất được giới hạn 30 sự kiện mỗi phút và 600 sự kiện mỗi lượt ghi danh. Ảnh được giới hạn 300 KB sau giải mã, kiểm tra dấu hiệu JPEG; tổng dung lượng ảnh mặc định 2 GiB, có thể cấu hình bằng biến môi trường.

Tên thư mục sinh viên được chuẩn hóa, tên ảnh có thành phần ngẫu nhiên. Bản ghi lưu đường dẫn ảnh, exam\_id và enrollment\_id; không lưu toàn bộ ảnh nhị phân vào SQLite. Tuyến đọc ảnh yêu cầu phiên giảng viên, không công khai thư mục snapshots.

Nhật ký cục bộ có thể tăng trước khi POST hoàn tất. Code chưa có hàng đợi gửi lại bền vững khi mất mạng, nên số cảnh báo phía sinh viên có thể khác số sự kiện lưu trên server. Các kiểm tra kích thước và kiểu tệp không thay thế đánh giá chất lượng bằng chứng.

\section{CSDL, sao lưu và phục hồi}
Đường dẫn DB được xác định theo thư mục backend, không theo thư mục làm việc của terminal. SQLite sử dụng WAL, synchronous=NORMAL, foreign\_keys=ON và busy\_timeout=5000; thời gian chờ kết nối là 5 giây \cite{sqlitewal}.

Khởi tạo kiểm tra toàn vẹn, tạo các bảng thiếu và chạy migrate\_db. Đây là migration tùy chỉnh theo cấu trúc hiện có, chưa phải hệ thống migration tổng quát có lịch sử phiên bản như Alembic. Backend có cơ chế sao lưu trước cập nhật schema khi cần và sao lưu định kỳ bằng SQLite backup API. Job sao lưu chạy mỗi sáu giờ; cơ chế lưu giữ giữ tối đa bảy bản sao định kỳ theo mẫu tên.

Ảnh được quét mỗi giờ và xóa khi captured\_at cũ hơn 24 giờ; đường dẫn được đặt rỗng, metadata còn lại. Thời hạn thực tế phụ thuộc lúc quét và thời gian dịch vụ hoạt động. captured\_at vẫn do client cung cấp nhưng đã được giới hạn độ lệch khi tiếp nhận.

recover\_db.py đọc nguồn và tạo DB phục hồi mới. Quá trình phục hồi không bảo đảm mọi dữ liệu trong trang hỏng đều còn nguyên. Các thao tác đổi tên/thay DB cần thực hiện khi backend đã dừng. Ngoài ra, handler DatabaseError trả 503 với thông báo thống nhất và rollback session khi có thể.

\section{Giao diện và triển khai đa thiết bị}
Không gian giảng viên gồm quản lý đề, yêu cầu vào thi và giám sát. Dashboard tải lại dữ liệu theo chu kỳ 10 giây; danh sách vi phạm hiện yêu cầu tối đa 500 dòng dù backend có limit/offset. Vì vậy thống kê tổng có thể lớn hơn số dòng đang hiển thị.

Web di động dùng video playsInline, gọi play() rõ ràng, giữ tỷ lệ ảnh và giảm chu kỳ AI xuống 500 ms. Khi chuyển ứng dụng, camera có thể ngừng cung cấp khung hình; hệ thống không tuyên bố quan sát liên tục trong nền. Khi có frontend/dist, FastAPI phục vụ giao diện cùng API và trả SPA fallback cho đường dẫn phù hợp.


Phần Flutter có camera, màn hình thi và các service HTTP. Tuy nhiên, hàm phân loại còn dựa trên chu kỳ thời gian thay vì mô hình AI; phần ước lượng độ sáng xử lý byte ảnh nén chưa đúng như pixel RGB. Service chưa hoàn chỉnh với header phiên thi đã duyệt. Do đó, phạm vi chức năng được đánh giá trong báo cáo tập trung vào backend và web, chưa xác nhận APK/IPA native hoàn chỉnh.
\begin{figure}[H]\centering\begin{tikzpicture}[node distance=0.9cm and 0.7cm]
\node[rb] (pc) {Máy tính\\Web React};
\node[rb,right=of pc] (phone) {Điện thoại\\Chrome / Safari / PWA};
\node[rb,below=of pc] (edge) {Cloudflare\\Địa chỉ HTTPS};
\node[rb,below=of edge] (tunnel) {cloudflared trên laptop\\Chuyển tiếp localhost:8000};
\node[rb,below=of tunnel] (server) {FastAPI\\API + frontend/dist};
\node[rb,right=of server] (store) {SQLite + snapshots\\Bản sao lưu DB};
\draw[ra] (pc)--(edge);\draw[ra] (phone)--(edge);\draw[ra] (edge)--(tunnel);\draw[ra] (tunnel)--(server);\draw[ra] (server)--(store);
\end{tikzpicture}

\caption{Mô hình triển khai web qua HTTPS tunnel.}\end{figure}
